%************************************************
\chapter{Energy-Based Models}\label{ch:ebms}
%************************************************

Energy-Based Models (EBMs) \cite{lecun2006tutorial} are a general class of models that capture dependencies by associating 
a scalar energy value to each configuration of variables, acting as a measure of compatibility.
\begin{itemize}
    \item \textbf{Training}: The general principle is to find an energy function that associates low energies to correct 
    (observed) configurations of variables, and higher energies to incorrect (unobserved) configurations.
    \item \textbf{Inference}: Making a decision consists of setting the values of observed variables and finding values 
    for the missing variables that minimize the energy, i.e., finding the values most compatible with the observation. 
    In the simplest scenario with two variables $X$ and $Y$, where the first is observed, the model will produce an 
    answer $Y^*$ that is the most compatible with $X$. This procedure is formally defined as:
    \begin{equation}
        Y^*=\text{argmin}_{Y\in \mathcal{Y}}E(X,Y).
    \end{equation} 
    When the size of the set $\mathcal{Y}$ is small, we can simply compute the energy for all possible values and pick 
    the smallest. However, we typically deal with complex, high-dimensional spaces where exhaustive search is impractical. 
    The solution is to apply an inference procedure to obtain an approximate result, which may or may not be the global 
    minimum of the energy function. For example, for a continuous variable $Y$ and a smooth, well-behaved energy function, 
    gradient-based optimization algorithms can be utilized.
\end{itemize}
The scenario described above covers prediction, classification, and decision-making, answering the question: "Which value 
of Y is most compatible with this X?". One advantage of these models is that we can use energy-based models to answer different 
types of questions. (I) Ranking, "Is Y1 or Y2 more compatible with this X?"; (II) Detection, "Is this value of Y compatible 
with X?"; and (III) Conditional density estimation, "What is the conditional probability distribution over $\mathcal{Y}$ 
given X?".  \\
\\
In deterministic settings, we only require the system to assign the lowest energy to the correct value, allowing us to ignore 
the absolute scale of the remaining large energies. However, if these outputs need to be rigorously compared, combined, or 
utilized to quantify uncertainty, uncalibrated energies are insufficient. We can calibrate energies by transforming them 
into a normalized probability distribution using the Gibbs distribution:
\begin{equation}
    p_{\theta}(\mathbf{x}) = \frac{e^{-E_{\theta}(\mathbf{x})}}{Z_{\theta}}, \qquad Z_{\theta}=\int_{\mathbf{x}\in\mathcal{X}} e^{-E_{\theta}(\mathbf{x})}d\mathbf{x}
\end{equation}
where the denominator $Z_{\theta}$ is called the \textit{partition function} ensuring normalization: 
\begin{equation}
    \int_{\mathbf{x}\in\mathcal{X}}p_{\theta}(\mathbf{x})d\mathbf{x} = 1.
\end{equation} 
Under this framework, lower energy values correspond to higher probabilities and the partition function $Z_{\theta}$ is a normalisation 
constant that ensures the probabilities sum to one. \\
\\
We can train probabilistic EBMs using Maximum Likelihood Estimation (MLE), defined by minimizing the negative log-likelihood 
of the empirical data:
\begin{align}
    \mathcal{L}(\theta) &= -\mathbb{E}_{\mathbf{x}\sim p_{data}}[\log p_{\theta}(\mathbf{x})] = -\mathbb{E}_{\mathbf{x}\sim p_{data}} \left[\log \frac{e^{-E_{\theta}(\mathbf{x})}}{Z_{\theta}}\right] \\
    & = \mathbb{E}_{\mathbf{x}\sim p_{data}} \left[E_{\theta}(\mathbf{x})\right] + \log Z_{\theta}
\end{align}
where $Z_{\theta}$ is the partition function defined in Equation --. The first term encourages the model to assign low energy to 
data points, while the second enforces normalization. The problem is that the partition function is intractable to compute for 
high-dimensional data, which makes the training of EBMs challenging. To overcome this, the EBM literature relies on two broad 
families of approximations: contrastive methods (such as Markov Chain Monte Carlo or Contrastive Divergence \cite{hinton2002training}), 
which approximate the model expectation, and non-contrastive/regularized methods (such as Score Matching \cite{hyvarinen2005estimation} or 
Denoising Autoencoders), which bypass the partition function entirely by restructuring the objective. \\
\\
In statistics, the score function of a probability distribution $p(\mathbf{x})$ is defined as the gradient of the log-probability 
density with respect to the input space $\mathbf{x}$: 
\begin{equation}
    \mathbf{s}(\mathbf{x}) = \nabla_{\mathbf{x}} \log p(\mathbf{x})
\end{equation}
When we apply this definition to the Gibbs distribution of our Energy-Based Model, a remarkable mathematical cancellation occurs:
\begin{align}
    \mathbf{s}{\theta}(\mathbf{x}) &= \nabla{\mathbf{x}} \log p_{\theta}(\mathbf{x}) \\
    &= \nabla_{\mathbf{x}} \left( -E_{\theta}(\mathbf{x}) - \log Z_{\theta} \right) \\
    &= -\nabla_{\mathbf{x}} E_{\theta}(\mathbf{x})
\end{align}
Because the partition function $Z_{\theta}$ is a constant with respect to the input $\mathbf{x}$, its gradient vanishes 
($\nabla_{\mathbf{x}} \log Z_{\theta} = 0$). Consequently, the score of an EBM is simply the negative gradient of its energy function. 
This allows us to evaluate the vector field pointing toward regions of higher probability without ever calculating the normalization 
constant.  

\section{Score Matching} 


\section{Denoising Score Matching} 






\begin{summarybox}[label={box:latent_math}]{EBMs and Score Matching}
    This box summarizes the main mathematical equations of EBMs and Score Matching:

    \medskip
    \textbf{1. Training Objective of VAEs:}
    \begin{equation}
        \mathcal{L}(\theta, \phi; \mathbf{x}) = \mathbb{E}_{z\sim q_{\phi}(\mathbf{z}|\mathbf{x})} \left[ \log p_{\theta}(\mathbf{x}|\mathbf{z}) \right] - D_{KL}(q_{\phi}(\mathbf{z}|\mathbf{x}) || p_{\theta}(\mathbf{z}))
    \end{equation}
    
    \medskip
    \textbf{2. Reparameterization Trick:}
    \begin{equation}
        \mathbf{x}_{\text{synthetic}} = G(\mathbf{z}_{\text{anatomy}} \oplus \mathbf{z}_{\text{pathology}} \mid \mathbf{y})
    \end{equation}
    
    %\vsub
    {\small \color{boxaccent}\textbf{Note:} Ensure latent dimensions $z_i$ and $z_j$ maintain near-zero mutual information to preserve explainability.}
\end{summarybox} 



%*****************************************
%*****************************************
%*****************************************
%*****************************************
%*****************************************